\chapter*{Özet}
\addcontentsline{toc}{chapter}{Özet}
Türkiye'den Schengen kısa süreli vizesine başvuranlar, birden fazla resmî sayfaya dağılmış ve
ülkeden ülkeye değişen gereksinimlerden yaklaşık on iki belgelik bir dosya hazırlar. Eksik bir
belge, süresi dolmak üzere olan bir pasaport ya da iki belgede farklı yazılmış bir isim, ödenen
ücrete ve randevuya mal olabilir. Genel amaçlı yapay zekâ asistanları bu sorulara akıcı cevaplar
verir; ancak bir gereksinimin hangi resmî cümleden geldiğini, hâlâ geçerli olup olmadığını ya da
başvuru sahibinin kendi belgelerinin birbiriyle tutarlı olup olmadığını göstermez.

\project{}, resmî gereksinimleri her belge ve kuralın kaynak cümlesini gösterdiği, sürümlenen
verilere dönüştürür. Başvuru sahibinin gideceği ülke ve durumu için belge listesi oluşturur,
yüklenen belgeleri şifreli saklar ve belgelerin temel bilgilerini deterministik bir kural motoruyla
kontrol eder: eksik belgeler, fazla eski belgeler, belgeler arasındaki ve seyahatle olan
uyuşmazlıklar. Her bulgu kuralı, ilgili belge ve alanları ve resmî kaynağı gösterir.

Güz dönemi prototipi Hollanda ve Norveç'i, çalışan ve öğrenci başvuru sahipleri için kapsar.
Etiketli 17 hata içeren iki sentetik başvuru dosyasında kural motoru her hatayı ve yalnızca onu
bulmuş, hatasız dosyalarda hiçbir bulgu vermemiştir. Belge görüntülerindeki bilgileri artık bir dil
modeli okumakta, başvuru sahibi bunları onaylamaktadır. Aynı görüntülerle aynı model tek başına 17
hatanın 16'sını bulmuş, ancak dosya başına ortalama 6,5 ek uyarı vermiştir; bunların çoğu resmî
kaynakların istemediği şeylerdir. Kural motoruyla birlikte 17 hatanın tamamını bulmuş ve hiç ek
uyarı vermemiştir.
Bahar döneminde çalışma, etiketli bir sentetik veri setine, insan değerlendiricilerle
karşılaştırmaya ve resmî sayfaların takibine genişletilecektir.

\noindent\textbf{Anahtar kelimeler:} vize başvurusu, belge doğrulama, kural motoru,
görüntü-dil modeli, kaynağa dayalı bilgi
